\chapter{Att kontrollera ett påstående}
\label{app:verifiera}

Det här materialet gör, som allt undervisningsmaterial, en rad
påståenden.  De flesta går att pröva vid datorn: att programmet kraschar
när filen saknas ser ni själva genom att köra det.  Men några av
påståendena handlar om vad forskningen, dokumentationen eller standarder
säger, och dem kan ingen körning avgöra:
\begin{enumerate}
  \item att nybörjare tar det som läses ur en fil för tal fastän det är
    strängar, och att materialets kontraster därför måste upplevas
    samtidigt för att göra skillnaden synlig (\cref{sec:egna,sec:csv});
  \item att primärminnet nås på nanosekunder och sekundärminnet på
    mikro- till millisekunder (\cref{sec:vad-ar-en-fil});
  \item att Pythons filfunktioner gör det materialet säger: att
    \mintinline{python}{with} stänger filen även vid fel, att
    \mintinline{python}{write} inte lägger till radbrytningar, att en
    slinga över filen är minnessnål, att text- och binärläge skiljer sig
    i avkodningen, och att en saknad fil ger
    \mintinline{python}{FileNotFoundError}
    (\cref{sec:oppna,sec:saknas,sec:lasa-skriva,sec:typer});
  \item att CSV är ett mycket vanligt men aldrig formellt standardiserat
    format, att modulen \mintinline{python}{csv} tar hand om kommatecken
    och citattecken i fält, och att JSON är en internetstandard som kan
    uttrycka nästlade strukturer (\cref{sec:csv,sec:json});
  \item att DRY, SRP och KISS är etablerade principer med de upphov
    materialet anger (\cref{sec:egna,sec:csv}).
\end{enumerate}
Hur vet man om sådant stämmer?  Samma fråga möter er varje gång ni läser
en lärobok, en webbsida eller ett svar från en språkmodell --- och varje
gång ni själva ska skriva en rapport.  Den här bilagan sammanfattar
arbetsgången; \cref{app:missuppfattningar,app:minne,app:dokumentation,%
app:format} visar den genomförd på de fyra första påståendena, och
\cref{tab:paastaaenden-files} sammanfattar frågorna och var svaren står.
Det femte påståendet är redan kontrollerat i andra föreläsningar och
görs inte om här.  Arbetsgången beskrivs utförligare i bilagan med samma
namn i föreläsningen \emph{Algoritmiskt tänkande}.

\section{Gör påståendet till en fråga}

Ett påstående går att kontrollera först när det är tydligt vad som skulle
kunna visa att det är \emph{fel}.  Formulera det därför som en fråga med
ett svar.  \enquote{Stänger \mintinline{python}{with} filen även när ett
fel inträffar?} kan besvaras med ja eller nej; \enquote{filhantering är
viktigt} kan det inte, eftersom \enquote{viktigt} är ett omdöme och inte
ett faktum.  Ofta döljer sig flera påståenden i ett: \enquote{CSV är ett
vanligt standardformat} är två frågor --- är det \emph{vanligt}, och är
det en \emph{standard} --- och i det här fallet får de olika svar.

\section{Välj rätt sorts källa}

Olika frågor besvaras av olika källor.
\begin{description}
  \item[Dokumentationen] för vad ett verktyg \emph{gör}.  Vad
    \mintinline{python}{open} gör med sitt läge avgörs av Pythons egen
    dokumentation, inte av en lärobok som återger den.
  \item[Standarden] för vad ett format \emph{är}.  För CSV och JSON är
    det de dokument (RFC:er) som beskriver formaten, och deras egen
    uppgift om vilken status de har.
  \item[Primärkällan] för var en idé \emph{kommer ifrån}, och för ett
    specifikt resultat.  Vill man veta vad studenter faktiskt gjorde
    läser man studien som mätte det.
  \item[Läroböcker] för storleksordningar och etablerade fakta, med
    förbehållet att siffrorna är avrundade och beror på maskinen.
  \item[Översikter] (\foreignlanguage{english}{reviews}) för om något är
    \emph{etablerat}, eller för vad forskningen \emph{sammantaget}
    säger.
\end{description}
Var man hittar dem: dokumentationen på \url{docs.python.org}, RFC:erna
hos \url{rfc-editor.org}, och för forskningen universitetsbibliotekets
söktjänst och databaser som Scopus, Web of Science, IEEE Xplore, DBLP
(datalogi) och OpenAlex (öppen).  I den här föreläsningen behövdes ingen
ny litteratursökning: forskningskällorna är återanvända ur kursens
sammanställning av nybörjares missuppfattningar, och de övriga
påståendena besvaras av dokument man vet finns och slår upp direkt ---
\emph{kända dokument}.

\section{Sök så att du kan ha fel}

Den som söker på namn hon redan känner till hittar det hon väntar sig.
Sök därför \emph{ämnesdrivet} --- på begrepp, inte författare --- och
sök också efter \emph{invändningar}.  För ett känt dokument gäller
motsvarigheten: läs också det som talar emot påståendet \emph{i samma
dokument}.  Både \cref{app:dokumentation,app:format} innehåller sådana
fall --- dokumentationen som rekommenderar \mintinline{python}{with}
säger också att standardkodningen beror på datorn, och RFC:n som
beskriver CSV säger i sin första mening att den inte är någon standard.
Ett påstående som överlevt en sådan läsning är värt mer än ett som bara
fått sin bekräftelse.

\section{Läs i fulltext, och skriv ner hur}

En träff i en databas är inte ett belägg, och en titel är det inte
heller.  Läs hela texten och leta upp \emph{det ställe} som styrker
påståendet --- ett ordagrant citat med sidnummer eller avsnitt.  Hittar
ni inget sådant ställe stöder källan inte påståendet.  Räkna aldrig en
källa ni inte läst: \cref{app:missuppfattningar} innehåller ett fall där
en källa som var läst i fulltext drogs tillbaka från webben dagen efter;
kontrollen gjordes i den sparade kopian, och det står uttryckligen så.

Skriv sedan ner hur ni gick till väga, så att någon annan kan följa
spåret.  I det här materialet står den anteckningen som en kommentar
ovanför varje post i litteraturlistans källfil, med fasta fält:
\begin{description}
  \item[\texttt{CLAIM}] vilket påstående källan ska styrka;
  \item[\texttt{FOUND-VIA}] hur den hittades --- sökfråga och databas,
    eller \enquote{känt dokument};
  \item[\texttt{PICKED}] varför just den, framför andra träffar;
  \item[\texttt{QUOTE}] det ordagranna citatet, med sida;
  \item[\texttt{VERIFIED}] att fulltexten lästs, och var;
  \item[\texttt{COUNTER}] motsökningen och vad den gav;
  \item[\texttt{DATE}] när.
\end{description}
Det tar ett par minuter per källa, och det är det enda som gör det
möjligt att om ett år svara på frågan \enquote{hur vet du det?}

\section{Dra slutsatsen --- med förbehåll}

Svaret skrivs ut tillsammans med det som talar emot.  Alla fyra
bilagorna svarar \enquote{ja, med förbehåll}
(\cref{tab:paastaaenden-files}): missuppfattningen är belagd, men en av
källorna är inte längre öppet tillgänglig; minnestiderna stämmer som
storleksordningar, inte som mätvärden för just er dator; dokumentationen
säger det materialet säger, men standardkodningen beror på datorn; och
CSV är mycket riktigt vanligt, men ingen standard --- det är JSON.
Förbehållen är inte svagheter i svaret.  De \emph{är} svaret, och de
avgör hur påståendet används i materialet: därför står det i
\cref{sec:csv} att RFC:n är informativ och att varianter förekommer, i
stället för att CSV kallas standard rakt av.

\begin{table}[htbp]
  \centering
  \caption{Påståendena i den här föreläsningen som frågor, och var
    svaret står.}
  \label{tab:paastaaenden-files}
  \begin{tabular}{@{}p{0.52\linewidth}l@{}}
    \toprule
    Fråga & Var svaret står \\
    \midrule
    Tar nybörjare det som läses ur en fil för tal fastän det är
      strängar? & \Cref{app:missuppfattningar} \\
    Måste en kontrast upplevas samtidigt för att kunna urskiljas? &
      \Cref{app:missuppfattningar} \\
    Nås primärminnet på nanosekunder och sekundärminnet på mikro- till
      millisekunder? & \Cref{app:minne} \\
    Gör Pythons filfunktioner det materialet säger? &
      \Cref{app:dokumentation} \\
    Är CSV vanligt, är det en standard, och tar modulen
      \mintinline{python}{csv} hand om specialfallen?  Är JSON en
      standard? & \Cref{app:format} \\
    Är det bra att inte upprepa sig i ett program, och kommer principen
      \foreignlanguage{english}{DRY} från Hunt och Thomas? &
      Redan belagt i \emph{Algoritmiskt tänkande}, bilaga E \\
    Kommer \foreignlanguage{english}{single responsibility} från Robert
      C. Martin, och \foreignlanguage{english}{KISS} från Kelly Johnson?
      Är de etablerad praxis? & Redan belagt i \emph{Funktioner},
      bilaga B och C \\
    \bottomrule
  \end{tabular}
\end{table}

\section{En anmärkning om språkmodeller}

I andra föreläsningars bilagor gav sökningarna hundratals träffar, och en
språkmodell användes för att \emph{sortera} dem.  Här gjordes ingen ny
sökning, så ingen sortering behövdes; men en språkmodell användes för
att hämta dokumenten, hitta citaten i dem och skriva anteckningarna.
Lägg märke till arbetsfördelningen: varje citat är hämtat ur källans
text, inte ur modellens minne, och varje källa ska läsas och godkännas av
en människa --- hur långt den granskningen kommit står i början av varje
bilaga.  Använd gärna språkmodeller för att hitta och sortera; låt dem
aldrig vara det sista ledet.


\chapter{Är missuppfattningen belagd?}
\label{app:missuppfattningar}
\chapterprecis{Författaren har ännu inte granskat resultaten i den här
  bilagan i sin helhet.}

Den här föreläsningen är på två ställen byggd kring en sak nybörjare
brukar missförstå: att det som läses ur en textfil är strängar, även när
det ser ut som tal.  I \cref{sec:egna} ställs frågan hur man får tillbaka
talen 3 och 9 ur raden \mintinline{text}|3,9| innan tolkningen visas,
och i \cref{sec:csv} står omvandlingen med \mintinline{python}{int} kvar
även när modulen \mintinline{python}{csv} tagit över uppdelningen.  Att
kontrasterna dessutom visas samtidigt --- de två läsprogrammen sida vid
sida, filens text bredvid det inlästa --- följer av en teori om lärande.
Frågan som ska besvaras är: \emph{är missuppfattningen belagd, och
håller skälet till att kontrasterna visas som de gör?}

\section{Metod}
\label{sec:miss-metod}

Ingen ny litteratursökning har gjorts för den här föreläsningen.
Källorna för missuppfattningen är hämtade ur en systematisk genomgång av
forskningen om nybörjares missuppfattningar som kursens lärare gjort
inom ett pågående forskningsarbete.  Den genomgången söktes ämnesdrivet
i flera databaser --- Semantic Scholar, OpenAlex, DBLP, arXiv, Web of
Science, IEEE Xplore och Scopus --- på frågan vilka missuppfattningar
nybörjare har i grundläggande programmering, och träffarna sorterades
först maskinellt och sedan för hand.  Källorna för lärandeteorin är de
som hela kursen använder i sina lärarnoter.

Att återanvända en källa är inte detsamma som att ha kontrollerat den.
För varje källa har därför tre saker gjorts om: påståendet har
formulerats så som \emph{den här} föreläsningen gör det; det ordagranna
citat föreläsningen stöder sig på har sökts upp igen i källan; och det
har antecknats vilken nivå kontrollen nådde --- fulltext, enbart
sammanfattning, eller inte alls.  Anteckningarna står i materialets
källfil för litteraturlistan, ovanför varje post.

Den här bilagan innehåller därför ingen träfflista och inga sökfrågor:
det vore att låtsas att en sökning gjorts som inte gjorts.  Vad som står
här är i stället vilka källor som bär vilket påstående, och hur långt
kontrollen av dem räcker.

% Author's note (verification and reproduction): no scholar session.
% Sources reused from ~/devel/edu/research/vt-prog-misconceptions
% (types.tex; misconceptions.bib) and from the computational-thinking
% deck's ltnotes.bib.  Re-verification 2026-09-03: Bosse2021Antipatterns
% re-read in full (arXiv 2104.12542, pdftotext), P_V5 confirmed verbatim;
% GuoMarkelZhang2020 NOT obtainable (author's copy withdrawn, ACM closed,
% no OA copy) -- quote carried on the paper's verification of 2026-06-29;
% Marton2015 pp. 45-47 as verified in the course, re-read 2026-09-04 on the
% user's reMarkable; MartonPang2006 abstract only (ERIC EJ733793) on
% 2026-09-03, full text read 2026-09-04 on the user's reMarkable (p. 199).
% Bib keys:
% Bosse2021Antipatterns, GuoMarkelZhang2020, Marton2015, MartonPang2006.

\section{Resultat}
\label{sec:miss-resultat}

\paragraph{Tas det som läses för tal?}  Ja.  Missuppfattningen är
katalogiserad som ett antimönster i verklig kod från
grundkursstudenter.  \Textcite{Bosse2021Antipatterns} byggde en katalog
över 41 antimönster ur 95 missuppfattningar i inlämnad kod, och
antimönstret \emph{missing type conversion} beskrivs så här:
\textcquote{Bosse2021Antipatterns}{Failure to convert, when necessary,
will make the data not be in the proper format so that the program can
work.  This error will be noticed during program execution}.  Att felet
gäller just data som läses från fil eller terminal står i den hittills
största samlingen av nybörjarförklaringar, drygt 16\,000 stycken
insamlade under tre år på en webbplats för programvisualisering:
\textcquote{GuoMarkelZhang2020}{When numerical data is read from files
or terminal input, they often start as strings.  If they are not
properly converted to numbers, it is still possible to use math
operators like + and > on them, which will perform string concatenation
and comparison}.  Det senare är precis vad frågan i \cref{sec:egna} är
byggd för att avslöja: den som stannar vid \mintinline{python}{split}
har ett program som ser ut att fungera men jämför strängar med tal.

\paragraph{Måste kontrasten upplevas samtidigt?}  Ja, enligt den teori
materialet är byggt på.  \Textcite[45--47]{Marton2015} visar att en aspekt bara
kan urskiljas om den varierar mot en invariant bakgrund, och att den
variationen måste upplevas samtidigt: \textcquote[46]{Marton2015}{if you cannot
see the greenness of any of three green, though otherwise different, things,
you cannot see what they have in common, either}.  Samma villkor formuleras i
den tidskriftsartikel materialets lärarnoter sedan tidigare hänvisade till:
\textcquote{MartonPang2006}{To discern an aspect, the learner must experience
potential alternatives, that is, variation in a dimension corresponding to that
aspect, against the background of invariance in other aspects of the same
object of learning}, och, uttryckligen om samtidigheten,
\textcquote[199]{MartonPang2006}{[a] quality X cannot be discerned without the
simultaneous experience of a mutually exclusive quality \(\sim\)X}.  Det är
därför de två läsprogrammen i \cref{sec:lasa-skriva} visas sida vid sida inför
valövningen, och därför körningen i \cref{sec:json} visar filens text och det
inlästa på samma bild.

\section{Diskussion och begränsningar}
\label{sec:miss-diskussion}

Tre saker begränsar svaret.

För det första är en av källorna inte längre tillgänglig.  Studien med
nybörjarförklaringarna kontrollerades i fulltext i författarens egen
kopia på webben; den kopian drogs tillbaka dagen efter att kursens andra
föreläsningar hämtat den, och förlagets version är stängd.  Citatet ovan
är läst om i den hämtade kopian, som kursen har sparat, men den som vill
kontrollera det själv får gå till förlaget.  Studien tillför att felet
gäller just data som läses från filer, vilket antimönsterkatalogen inte
säger ut.

För det andra var den lärandeteoretiska artikeln först läst enbart i
sammanfattning, eftersom den är stängd hos förlaget; den lästes sedan i
fulltext i författarens eget exemplar.  Den ställer samtidighetskravet
själv, som en av fyra nödvändiga betingelser för perceptuellt lärande,
och hänvisar för det till Marton och Tsui (2004); det skarpare
påståendet bärs därför av både artikeln och boken.

För det tredje säger ingen av källorna något om hur ofta felet
förekommer i just den här kursen.  Antimönsterkatalogen bygger på kod från
en grundkurs i ett annat land, delvis i C; förklaringarna kommer från
användare av en webbplats, inte från en kurs alls.  Att felet känns igen
här är inte samma sak som att frekvenserna överförs.  Kursens egna
diagnostiska frågor finns till för det.

\section{Slutsats}
\label{sec:miss-slutsats}

Ja, missuppfattningen är belagd, med förbehållet att en av de två källorna inte
längre är öppet tillgänglig och att ingen av dem mäter frekvensen i den här
kursens population.  Att data som läses in inte omvandlas till tal är ett
katalogiserat antimönster i nybörjarkod\autocite{Bosse2021Antipatterns}, och
att felet gäller just det som läses från filer och terminal står i studenters
egna förklaringar\autocite{GuoMarkelZhang2020}.  Att kontrasterna visas
samtidigt följer av variationsteorins villkor för
urskiljning\autocite[45--47]{Marton2015}, som artikeln materialet tidigare
hänvisade till formulerar på samma sätt\autocite[199]{MartonPang2006}.
Beläggen används till det de räcker till: att välja \emph{vilken} fråga som
ställs innan tolkningen visas, inte till att säga åt studenterna vad de tror.


\chapter{Är primärminnet så mycket snabbare?}
\label{app:minne}
\chapterprecis{Författaren har ännu inte granskat resultaten i den här
  bilagan i sin helhet.}

I \cref{sec:vad-ar-en-fil} sägs att primärminnet nås i storleksordningen
nanosekunder och sekundärminnet i mikro- till millisekunder.  Det är ett
faktapåstående om maskinvara, och det motiverar hela avvägningen mellan
de två minnena.  Frågan som ska besvaras är: \emph{stämmer
storleksordningarna?}

\section{Metod}
\label{sec:minne-metod}

Det här är en fråga om etablerade fakta, så källan är en lärobok, hämtad
som känt dokument: \emph{Operating Systems: Three Easy Pieces}, en fritt
publicerad lärobok i operativsystem vars kapitel läses direkt på
webben\autocite{ArpaciDusseau2025OSTEP}.  Den valdes framför de
klassiska läroböckerna i datorarkitektur därför att var och en kan
öppna den och kontrollera citaten; de klassiska böckernas tabeller gick
inte att läsa i fulltext.  Tre kapitel lästes: det om
sidersättningspolicyer, där boken räknar på kostnaden för minnes- och
diskåtkomst, och de om hårddiskar och SSD:er, där komponenttiderna
anges.

% Author's note (verification and reproduction): no scholar session.
% Chapters fetched 2026-09-03 from https://pages.cs.wisc.edu/~remzi/OSTEP/
% (vm-beyondphys-policy.pdf, file-disks.pdf, file-ssd.pdf), read with
% pdftotext; the passages quoted below are in sec. 22.1, 37.4 and 44.3.
% Bib key: ArpaciDusseau2025OSTEP.

\section{Resultat}
\label{sec:minne-resultat}

Boken använder just de storleksordningar materialet anger.  I
räkneexemplet om genomsnittlig åtkomsttid skriver författarna:
\textcquote[kap.~22.1]{ArpaciDusseau2025OSTEP}{Assuming the cost of
accessing memory (\(T_M\)) is around 100 nanoseconds, and the cost of
accessing disk (\(T_D\)) is about 10 milliseconds}.  Disktiden består av
komponenter som kapitlet om hårddiskar anger var för sig, bland dem
\textcquote[kap.~37.4]{ArpaciDusseau2025OSTEP}{The average seek time (4
milliseconds)}.  För SSD:er är läsningen av en sida
\textcquote[kap.~44.3]{ArpaciDusseau2025OSTEP}{typically quite fast, 10s
of microseconds or so}.  Primärminne på hundra nanosekunder,
sekundärminne på tiotals mikrosekunder till millisekunder: det är
materialets påstående.

\section{Diskussion och begränsningar}
\label{sec:minne-diskussion}

Siffrorna är läroboksvärden, inte mätningar.  Bokens minnes- och
disktider införs som \emph{antaganden} i ett räkneexempel, och
komponenttiderna anges som typiska värden som beror på tillverkare och
generation.  En modern SSD kan vara snabbare än tiotals mikrosekunder
och en gammal hårddisk långsammare än några millisekunder.  Materialet
säger därför storleksordningar, inte tal, och det är på den nivån
påståendet är kontrollerat.

\section{Slutsats}
\label{sec:minne-slutsats}

Ja, storleksordningarna stämmer: hundra nanosekunder för primärminnet
mot tiotals mikrosekunder (SSD) till millisekunder (hårddisk) för
sekundärminnet\autocite[kap.~22, 37 och 44]{ArpaciDusseau2025OSTEP},
med förbehållet att de är avrundade läroboksvärden som varierar med
maskinen.


\chapter{Gör Pythons filfunktioner det materialet säger?}
\label{app:dokumentation}
\chapterprecis{Författaren har ännu inte granskat resultaten i den här
  bilagan i sin helhet.}

Materialet gör en rad påståenden om vad Pythons filfunktioner gör.  I
\cref{sec:oppna} sägs att filen måste stängas för att det skrivna säkert
ska nå disken, och att \mintinline{python}{with} stänger den även när ett
särfall inträffar.  I \cref{sec:saknas} sägs att en saknad fil ger
\mintinline{python}{FileNotFoundError}.  I \cref{sec:lasa-skriva} sägs
att \mintinline{python}{write} inte lägger till radbrytningar, att varje
rad som läses behåller sin radbrytning, att en slinga över filen är
minnessnål, att \mintinline{python}{read} läser hela filen och att allt
som läses är strängar.  I \cref{sec:typer} sägs att textläget avkodar
filens bytes med en teckenkodning medan binärläget ger råa bytes, och
att standardkodningen beror på datorn.  Frågan som ska besvaras är:
\emph{stämmer allt detta?}

\section{Metod}
\label{sec:dok-metod}

Det här är frågor om vad ett verktyg gör, så den avgörande källan är
Pythons egen dokumentation, hämtad som känt dokument den 3 september
2026: handledningens avsnitt om att läsa och skriva
filer\autocite{PythonTutorialFiles}, referensen för
\mintinline{python}{open}\autocite{PythonOpen} och referensen för de
inbyggda särfallen\autocite{PythonExceptions}.  Sidorna hämtades som
text och varje påstående söktes upp i dem ordagrant.  Ingen
litteratursökning gjordes: dokumentationen är normerande för vad
funktionerna gör, och en studie som säger något annat skulle beskriva
ett fel i Python, inte i materialet.

% Author's note (verification and reproduction): no scholar session.
% Pages fetched with curl 2026-09-03 (tutorial/inputoutput.html,
% library/functions.html, library/exceptions.html), tags stripped, quotes
% located with a small Python script (quotes.py in the agent's scratch).
% Bib keys: PythonTutorialFiles, PythonOpen, PythonExceptions.

\section{Resultat}
\label{sec:dok-resultat}

\paragraph{Stängning och \mintinline{python}{with}.}  Ja.  Handledningen
säger att \textcquote{PythonTutorialFiles}{It is good practice to use the
with keyword when dealing with file objects.  The advantage is that the
file is properly closed after its suite finishes, even if an exception
is raised at some point}, och om den som inte gör det:
\textcquote{PythonTutorialFiles}{If you're not using the with keyword,
then you should call f.close() to close the file and immediately free up
any system resources used by it}.  Att det skrivna annars kan utebli står
i en varning: \textcquote{PythonTutorialFiles}{Calling f.write() without
using the with keyword or calling f.close() might result in the arguments
of f.write() not being completely written to the disk, even if the
program exits successfully}.

\paragraph{En saknad fil.}  Ja.  Referensen definierar
\mintinline{python}{FileNotFoundError} som det särfall som kastas
\textcquote{PythonExceptions}{when a file or directory is requested but
doesn't exist}.  Samma sida förklarar också varför materialets tidigare
exempelprogram, som fångade \mintinline{python}{EnvironmentError},
ändrats: det namnet är sedan Python 3.3 bara ett alias för
\mintinline{python}{OSError}, kvar \textcquote{PythonExceptions}{for
compatibility with previous versions}.

\paragraph{Läsa och skriva.}  Ja, på alla fyra punkterna.
\mintinline{python}{write} lägger inte till något:
\textcquote{PythonTutorialFiles}{f.write(string) writes the contents of
string to the file, returning the number of characters written}.  Raden
behåller sin radbrytning: \textcquote{PythonTutorialFiles}{a newline
character (\textbackslash{}n) is left at the end of the string, and is
only omitted on the last line of the file if the file doesn't end in a
newline}.  Slingan är minnessnål: \textcquote{PythonTutorialFiles}{For
reading lines from a file, you can loop over the file object.  This is
memory efficient, fast, and leads to simple code}.  Och
\mintinline{python}{read} utan argument läser allt --- med
dokumentationens egen brasklapp om att det är vårt problem om filen är
dubbelt så stor som datorns minne.  Att det lästa är strängar står i
samma avsnitt, om \mintinline{python}{f.read(size)}:
\textcquote{PythonTutorialFiles}{returns it as a string (in text mode)
or bytes object (in binary mode)}.

\paragraph{Text och binärt.}  Ja.  Referensen för
\mintinline{python}{open} säger att \textcquote{PythonOpen}{Files opened
in binary mode (including 'b' in the mode argument) return contents as
bytes objects without any decoding.  In text mode (the default, or when
't' is included in the mode argument), the contents of the file are
returned as str, the bytes having been first decoded using a
platform-dependent encoding or using the specified encoding if given},
och om lägena att \mintinline{python}{"w"} öppnar
\textcquote{PythonOpen}{for writing, truncating the file first} medan
\mintinline{python}{"a"} öppnar \textcquote{PythonOpen}{for writing,
appending to the end of file if it exists}.

\paragraph{Det som talar emot.}  Ett förbehåll kommer ur samma sidor.
Materialet sade tidigare att textläget \enquote{oftast} använder UTF-8;
det gör det inte av sig självt.  Referensen säger att
\textcquote{PythonOpen}{The default encoding is platform dependent
(whatever locale.getencoding() returns)}, och handledningen att
\textcquote{PythonTutorialFiles}{Because UTF-8 is the modern de-facto
standard, encoding="utf-8" is recommended unless you know that you need
to use a different encoding}.  Materialet säger därför numera i
\cref{sec:typer} att kodningen beror på datorn och att dokumentationen
rekommenderar att ange den.

\section{Diskussion och begränsningar}
\label{sec:dok-diskussion}

Svaret vilar helt på dokumentationen, som beskriver den aktuella
utgåvan av Python och kan ändras; datumet står i litteraturlistan.  De
påståenden materialet gör har varit stabila genom Python 3, men
standardkodningen är ett exempel på något som är på väg att ändras ---
dokumentationens rekommendation att alltid ange UTF-8 är just ett tecken
på det.  Materialets exempelprogram anger ingen kodning, av enkelhet;
det är ett medvetet val som bilagan här redovisar.

\section{Slutsats}
\label{sec:dok-slutsats}

Ja, filfunktionerna gör det materialet säger, med förbehållet att
standardkodningen i textläge beror på datorn och inte alltid är UTF-8.
Att \mintinline{python}{with} stänger filen även vid fel, att stängningen
behövs för att det skrivna ska nå disken, att \mintinline{python}{write}
inte lägger till radbrytningar, att lästa rader behåller sina, att slingan
över filen är minnessnål och att det lästa är
strängar\autocite{PythonTutorialFiles}; att en saknad fil ger
\mintinline{python}{FileNotFoundError}\autocite{PythonExceptions}; och att
text- och binärläge skiljer sig i avkodningen\autocite{PythonOpen} ---
allt står i dokumentationen med de ord materialet använder.


\chapter{Är CSV och JSON standarder?}
\label{app:format}
\chapterprecis{Författaren har ännu inte granskat resultaten i den här
  bilagan i sin helhet.}

I \cref{sec:csv} kallas CSV ett av de vanligaste formaten för
tabelldata, sägs vara beskrivet men inte standardiserat i RFC~4180, och
modulen \mintinline{python}{csv} sägs ta hand om kommatecken och
citattecken inuti fält.  I \cref{sec:json} kallas JSON en
internetstandard som kan uttrycka nästlade strukturer.  Det är flera
påståenden i ett, och de kräver olika slags svar.  Frågan som ska
besvaras är: \emph{är CSV vanligt, är det en standard, tar modulen hand
om specialfallen --- och är JSON en standard?}

\section{Metod}
\label{sec:format-metod}

Formaten beskrivs av kända dokument, som hämtades direkt den 3 september
2026: RFC~4180 för CSV\autocite{RFC4180}, RFC~8259 för
JSON\autocite{RFC8259}, samt Pythons dokumentation för modulerna
\mintinline{python}{csv}\autocite{PythonCSV} och
\mintinline{python}{json}\autocite{PythonJSON}.  Frågan om
\emph{status} besvaras av dokumenten själva: varje RFC anger i sin
ingress vilken kategori den tillhör.  Frågan om \emph{utbredning} är
svårare, för den är empirisk; här användes dokumentationens och RFC:ns
egna utsagor, och begränsningen i det redovisas nedan.  Ingen
litteratursökning gjordes.

% Author's note (verification and reproduction): no scholar session.
% RFCs fetched as plain text from rfc-editor.org and the two
% docs.python.org pages with curl, 2026-09-03; quotes located verbatim.
% Bib keys: RFC4180, RFC8259, PythonCSV, PythonJSON.

\section{Resultat}
\label{sec:format-resultat}

\paragraph{Är CSV vanligt?}  Ja, enligt de källor som beskriver det.
Pythons dokumentation inleder med att \textcquote{PythonCSV}{The
so-called CSV (Comma Separated Values) format is the most common import
and export format for spreadsheets and databases}, och RFC:n att formatet
\textcquote{RFC4180}{has been used for exchanging and converting data
between various spreadsheet programs for quite some time.  Surprisingly,
while this format is very common, it has never been formally
documented}.

\paragraph{Är CSV en standard?}  Nej.  RFC~4180 säger det själv i
ingressen: \textcquote{RFC4180}{This memo provides information for the
Internet community.  It does not specify an Internet standard of any
kind}.  Dokumentationen drar konsekvensen: \textcquote{PythonCSV}{The
lack of a well-defined standard means that subtle differences often
exist in the data produced and consumed by different applications}.  Det
är därför materialet säger \enquote{beskrivs} om RFC:n och inte
\enquote{standardiseras}.

\paragraph{Tar modulen hand om specialfallen?}  Ja, och de är just de
som RFC:n beskriver: \textcquote{RFC4180}{Fields containing line breaks
(CRLF), double quotes, and commas should be enclosed in double-quotes},
och ett citattecken inuti ett sådant fält dubbleras.  Dokumentationen
beskriver modulens uppgift som att \textcquote{PythonCSV}{write a single
module which can efficiently manipulate such data, hiding the details of
reading and writing the data from the programmer}, och föreskriver att
filen öppnas med \mintinline{python}{newline=""}:
\textcquote{PythonCSV}{If csvfile is a file object, it should be opened
with newline=''}.

\paragraph{Är JSON en standard?}  Ja.  RFC~8259 är, till skillnad från
RFC~4180, ett dokument på standardspåret --- \textcquote{RFC8259}{This is
an Internet Standards Track document} --- och beskriver
\textcquote{RFC8259}{a lightweight, text-based, language-independent
data interchange format} som \textcquote{RFC8259}{can represent four
primitive types (strings, numbers, booleans, and null) and two
structured types (objects and arrays)}.  Pythons modul implementerar
formatet enligt den föregående utgåvan av samma
specifikation\autocite{PythonJSON}.

\section{Diskussion och begränsningar}
\label{sec:format-diskussion}

Att CSV är \emph{vanligt} vilar på två utsagor från dem som beskriver
formatet, inte på en mätning.  Ingen studie av hur vanliga olika
filformat är har sökts, och påståendet står därför i materialet med
källan utsatt --- \enquote{Pythons dokumentation kallar det} --- i
stället för som ett eget faktum.  Frågan om status är däremot
avgjord av dokumenten själva, och svaret vände på materialets tidigare
formulering: CSV kallades ett standardformat, och det är JSON som är
det.

\section{Slutsats}
\label{sec:format-slutsats}

Ja, med ett viktigt förbehåll: CSV är ett mycket vanligt format enligt
både Pythons dokumentation\autocite{PythonCSV} och RFC:n som beskriver
det\autocite{RFC4180}, men RFC:n är informativ och formatet är ingen
standard; modulen \mintinline{python}{csv} hanterar just de specialfall
RFC:n beskriver --- kommatecken, citattecken och radbrytningar i fält
--- när filen öppnas som dokumentationen föreskriver.  JSON är däremot
en internetstandard på standardspåret\autocite{RFC8259}, med de fyra
enkla och två sammansatta typer materialet räknar upp.  Materialet
säger därför \enquote{standardformat} om JSON och \enquote{beskrivet i
RFC~4180} om CSV.
